\section{Π/Σ-types and the calculus of constructions}\label{sec:02-terminology-and-coc:02-pi-sigma-and-coc:1}

Section 3 built Π-types concretely, from \texttt{Fin n} and \texttt{Vec.replicate}, up to the general pattern \(\prod_{x:A} B(x)\). This section adds the second half of the picture, \textbf{Σ-types}, the dependent pair dual to the Π-type. It then zooms out to name the whole formal system these pieces belong to, the \textbf{calculus of constructions} (CoC), the core type theory underlying Lean. Universes get their \emph{own} formal typing rule in \hyperref[sec:06-rigor-check:03-typing-rules-and-safety:1]{Chapter 6, Section 3}, once Chapter 6 has built up the informal picture of \texttt{Type}/\texttt{Type 1} those rules formalize. Nothing here depends on having read that yet.

\subsection{Π-types, briefly recalled}\label{sec:02-terminology-and-coc:02-pi-sigma-and-coc:2}

\[
\prod_{x : A} B(x)
\]

This reads as ``a function that, given \(x : A\), returns a term of type \(B(x)\), a type allowed to mention \(x\).'' The \texttt{Vec.replicate} from Section 3, with type \texttt{(n : Nat) → Vec α n}, literally \emph{is} \(\prod_{n:\mathtt{Nat}}
\mathrm{Vec}\,\alpha\,n\). The surface syntax \texttt{(x : A) → B x} in Lean \textbf{is} \(\prod_{x:A} B(x)\), and \texttt{∀} is the same construct specialized to \texttt{B : A → Prop}, as Chapter 4 makes precise. Every \texttt{∀ n : Nat, P n} written from that point on is already a Π-type in exactly this sense.

\begin{progcorner}

Even the \texttt{TypeVar} construct in Python (its own fix for generic functions like \texttt{identity}) cannot express the type of \texttt{Vec.replicate}:

\end{progcorner}

\begin{lstlisting}[language=Python,style=python]
from typing import TypeVar
T = TypeVar("T")

def identity(x: T) -> T:      # generic over a TYPE (*\textemdash{}*) fine, this is what TypeVar is for
    return x

# There is no Python type-hint equivalent of:
#   def replicate(a: T, n: <this specific Nat>) -> Vec[T, <that same Nat>]
# because "the type mentions the VALUE of n" is not expressible by any
# combination of ordinary hints or TypeVar.
\end{lstlisting}

This is the honest answer to why Lean needs a whole extra concept here. It is not that the type hints in Python are missing a minor convenience. It is that \emph{no} mainstream static type system, not in Python, Java, C\#, nor TypeScript, has a construct for ``the type mentions a specific runtime value,'' because none of them needed the level of precision a proof assistant requires. Π-types are precisely that missing construct, made rigorous.

\textbf{A second worked example, where the codomain is a genuinely different type, not just a differently-sized version of the same type.} \texttt{Vec.replicate} always returns \emph{some} \texttt{Vec}. The following function returns a \texttt{Nat} or a \texttt{Bool} depending on which was asked for, an even sharper case of ``the type mentions the value of the argument'':

\begin{lstlisting}
def pick (b : Bool) : (if b then Nat else Bool) :=
  match b with
  | true => (42 : Nat)
  | false => (true : Bool)

#check @pick   -- pick : (b : Bool) (*$\rightarrow$*) if b = true then Nat else Bool
#eval pick true   -- 42
#eval pick false  -- true
\end{lstlisting}

\textbf{A counterexample.} Forcing the result of \texttt{pick} to a single fixed type (dropping the dependency) is rejected outright, because the two branches genuinely disagree on type. There is no ordinary, non-dependent function type that could describe \texttt{pick}:

\begin{lstlisting}
def badPick (b : Bool) : Nat :=
  match b with
  | true => 42
  | false => true
\end{lstlisting}

\begin{lstlisting}
error: Type mismatch
  true
has type
  Bool
but is expected to have type
  Nat
\end{lstlisting}

\subsection{\texorpdfstring{\texttt{Prop} as a special, proof-irrelevant universe}{Prop as a special, proof-irrelevant universe}}\label{sec:02-terminology-and-coc:02-pi-sigma-and-coc:3}

The \texttt{Prop} type in Lean (Chapter 4) sits alongside \texttt{Type} as its own universe, with one extra rule, \textbf{proof irrelevance}. Any two terms of the same type \texttt{P : Prop} are considered definitionally equal, since a proof carries no computational content beyond the bare fact that \emph{a} proof exists:

\begin{lstlisting}
theorem two_proofs (h1 h2 : 2 + 2 = 4) : h1 = h2 := rfl
\end{lstlisting}

\texttt{rfl} succeeds no matter \emph{how} \texttt{h1} and \texttt{h2} were each proved, whether by \texttt{rfl} directly, by a long tactic block, or by an entirely different chain of lemmas, because Lean does not distinguish between different proofs of the same \texttt{Prop}. This is what makes the slogan of Curry--Howard, ``propositions are types, proofs are terms'' more than a slogan. \texttt{P : Prop} really is a type, \texttt{h : P} really is an ordinary term of that type built by ordinary \(\lambda\)/application/Π-type machinery, and the \emph{only} difference from an ordinary \texttt{Type} is that Lean does not care \emph{which} term of \texttt{P} was produced, only that one was produced.

\subsection{Σ-types, the dependent pair, dual to the Π-type}\label{sec:02-terminology-and-coc:02-pi-sigma-and-coc:4}

Where the Π-type generalizes the ordinary function type \(\to\), a \textbf{Σ-type} (dependent pair / dependent sum) generalizes \(\times\) (the ordinary product):

\[
\sum_{x:A} B(x)
\]

This is a pair \(\langle a, b\rangle\) with \(a : A\) and \(b : B(a)\), where the type of the \emph{second} component is allowed to depend on the \emph{value} of the \emph{first} component.

\textbf{Why ``sum,'' if it generalizes ×?} The name is not a mismatch between the Σ-type and the product. It names the more fundamental view underneath both. \(\sum_{x:A} B(x)\) is, categorically, a disjoint union (coproduct) of the family \(B\) \emph{indexed by} \(A\), one tagged copy of \(B(x)\) for every \(x \in A\), glued together as separate parts that never overlap. The Curry--Howard table in \hyperref[sec:04-propositions-and-proofs:01-prop:1]{Chapter 4, Section 1} will list \texttt{∧} as a \textbf{product type} and \texttt{∨} as a \textbf{sum (coproduct) type}, side by side with \texttt{∃} as a \textbf{Σ-type}, as though all three were unrelated rows. They are not. \texttt{∧} and \texttt{∨} are both special cases of this same Σ-type construction. Which connective you get depends on which of the two ingredients, the index type \(A\) or the family \(B\), is held fixed and which is allowed to vary.

\begin{itemize}
\tightlist
\item
  \textbf{Constant family, varying index.} If \(B(x)\) is the \emph{same} type \(B\) for every \(x\), the disjoint union of \(|A|\) tagged copies of \(B\) coincides with the ordinary product \(A \times B\) (a value of \(A\), tagging \emph{which} copy, paired with a value of \(B\)). This is the ``dependent pair'' reading used throughout this section, and it is exactly what \texttt{∧} in Chapter 4 will turn out to be. \(P \wedge Q\) is \(\sum_{\_ : P} Q\), an index type \(P\) with a family constantly \(Q\), so the ``which copy'' tag is itself a \emph{proof} rather than a value.
\item
  \textbf{Two-point index, varying family.} If \(A\) is a two-element type (the \texttt{Bool} type in Lean, or ``which side'' of an \texttt{Or}) and \(B(\mathrm{true}) = P\), \(B(\mathrm{false}) = Q\) for two unrelated propositions, the \emph{same} construction \(\sum_{x:\mathrm{Bool}} B(x)\) collapses instead to \(P \sqcup Q\), the ordinary sum/coproduct type that \texttt{P ∨ Q} in Chapter 4 will turn out to be, built from \texttt{Or.inl}/\texttt{Or.inr}.
\end{itemize}

The Σ-type is called a \emph{sum} because it always literally is an indexed disjoint union. The product reading (\texttt{×}, and \texttt{∧} as its \texttt{Prop} special case) is just the constant-family special case of that same construction, not a second, unrelated thing.

The same split applies to the Π-type, the other way round. \(\prod_{x:A} B(x)\) is always literally an indexed \emph{product}. For a constant family, this collapses to the ordinary function type \(A \to B\) (an exponential, written \(B^A\)), not to \(A \times B\). This is why the Π-type is never confused with \texttt{×} the way the Σ-type is with \texttt{∨}, and no equivalent puzzle arises on that side.

This shape has already appeared, without the name. The \texttt{Fin n} from Section 3 is, under the hood, exactly this pair:

\begin{lstlisting}
#print Fin
-- structure Fin (n : Nat) : Type
-- fields:
--   Fin.val  : Nat
--   Fin.isLt : ↑self < n
\end{lstlisting}

It is a \texttt{Nat} value \texttt{val}, paired with a proof \texttt{isLt} whose \emph{statement} (\texttt{↑self < n}) mentions the first field itself (\texttt{self.val}, printed with the \texttt{↑} coercion arrow). That is \(\Sigma\), concretely. The type of the second field depends on the value of the first. Here is the general construct, spelled out with the actual \texttt{Sigma} type in Lean, pairing a bound with an actual number below it:

\begin{lstlisting}
def mySigma : (*$\Sigma$*) n : Nat, Fin n := (*$\langle$*)3, (*$\langle$*)2, by decide(*$\rangle$*)(*$\rangle$*)
#eval mySigma.fst        -- 3
#eval mySigma.snd.val    -- 2
\end{lstlisting}

\texttt{mySigma.fst} extracts the first component with no fuss at all. It is ordinary data, just like extracting \texttt{.1} from an ordinary pair. This extractability matters, because it is exactly what \texttt{∃} (next) does \emph{not} allow, and the reason why is the most subtle point in this section.

\textbf{A second worked example, mirroring the Π-type one above.} Just as the codomain of a Π-type can be a genuinely different type per argument, the second component of a Σ-type can be a genuinely different type per first component, not just a differently-sized version of one fixed type:

\begin{lstlisting}
def mySigma2 : (*$\Sigma$*) b : Bool, if b then Nat else String :=
  (*$\langle$*)true, (42 : Nat)(*$\rangle$*)
def mySigma3 : (*$\Sigma$*) b : Bool, if b then Nat else String :=
  (*$\langle$*)false, "hi"(*$\rangle$*)

#eval mySigma2.fst  -- true
#eval mySigma2.snd  -- 42
#eval mySigma3.fst  -- false
#eval mySigma3.snd  -- "hi"
\end{lstlisting}

\textbf{A counterexample.} The type of the second component is dictated by the \emph{specific value} supplied as the first component, not freely choosable alongside it. Supplying a \texttt{String} alongside \texttt{true} (whose dependent type here is \texttt{Nat}) is rejected:

\begin{lstlisting}
def badSigma : (*$\Sigma$*) b : Bool, if b then Nat else String :=
  (*$\langle$*)true, "oops"(*$\rangle$*)
\end{lstlisting}

\begin{lstlisting}
error: Application type mismatch: The argument
  "oops"
has type
  String
but is expected to have type
  if true = true then Nat else String
in the application
  (*$\langle$*)true, "oops"(*$\rangle$*)
\end{lstlisting}

This dependent-pair shape is also the ``structure bundling data + proofs'' pattern that recurs throughout this book. The \texttt{⟨op, id, inv, assoc, ...⟩} of \texttt{Group G} is, underneath the \texttt{structure} sugar in Lean, an iterated Σ-type, a witness \texttt{op}, paired with a witness \texttt{id}, paired with \texttt{assoc}, whose \emph{type} depends on the values of \texttt{op} and \texttt{id} supplied earlier in the same structure. This dependent-pair shape is exactly what lets Chapter 3 say that structures can bundle proofs alongside data. Every such structure is, underneath, silently an appeal to the Σ-type.

\textbf{A caveat about \texttt{∃}, worth being precise about.} Chapter 4 will read \texttt{∃ x : α, P x} as ``a structure, a witness value plus a proof that the witness satisfies \texttt{P},'' in effect a Σ-type, with \texttt{P : α → Prop} playing the role of \texttt{B}. The actual \texttt{Exists} type in Lean is, however, \emph{not literally} the Σ-type \texttt{Sigma}. \texttt{Exists} is specifically built to land in \texttt{Prop}, and by proof irrelevance (above), that means the witness cannot be \emph{extracted} from an \texttt{∃}-proof computationally:

\begin{lstlisting}
def bad (h : (*$\exists$*) n : Nat, n > 0) : Nat := h.1
\end{lstlisting}

\begin{lstlisting}
error(lean.projNonPropFromProp): Invalid projection: Cannot project a
value of non-propositional type
  Nat
from the expression
  h
which has propositional type
  (*$\exists$*) n, n > 0
\end{lstlisting}

Compare this directly to \texttt{mySigma.fst} two paragraphs up, which worked with no error at all. Same pair-shape, different universe, and that difference is exactly why one extracts and the other does not. If \texttt{∃} \emph{did} allow extracting a witness, two different (but both equally valid) choices of witness could produce different results from a ``proof-irrelevant'' input, contradicting proof irrelevance itself. \texttt{Sigma} (the actual, \texttt{Type}-valued dependent pair, matching the bundled data of a \texttt{structure}) \emph{does} support projecting out its first component, precisely because it does not carry the irrelevance guarantee of \texttt{Exists}. In short, \texttt{∃} has the same shape as the Σ-type but lives in a universe where the witness is unobservable. That is what makes it a \emph{restricted} cousin of the Σ-type rather than literally the Σ-type. The \texttt{structure}-based bundling this book uses throughout for \texttt{Group}/\texttt{Ring}/\texttt{Module} genuinely is \texttt{Sigma}-like (extractable), while an \texttt{∃}-statement genuinely is not.

\subsection{Π-types and Σ-types, side by side}\label{sec:02-terminology-and-coc:02-pi-sigma-and-coc:5}

Everything above in one table, gathering the two constructs' notation, readings, and special cases for direct comparison.

\begin{longtable}[]{@{}
  >{\raggedright\arraybackslash}p{(\linewidth - 4\tabcolsep) * \real{0.3333}}
  >{\raggedright\arraybackslash}p{(\linewidth - 4\tabcolsep) * \real{0.3333}}
  >{\raggedright\arraybackslash}p{(\linewidth - 4\tabcolsep) * \real{0.3333}}@{}}
\toprule\noalign{}
\begin{minipage}[b]{\linewidth}\raggedright
Aspect
\end{minipage} & \begin{minipage}[b]{\linewidth}\raggedright
Π-type \(\prod_{x:A} B(x)\)
\end{minipage} & \begin{minipage}[b]{\linewidth}\raggedright
Σ-type \(\sum_{x:A} B(x)\)
\end{minipage} \\
\midrule\noalign{}
\endhead
\bottomrule\noalign{}
\endlastfoot
Read as & ``for every \(x:A\), a term of \(B(x)\)'' & ``some \(x:A\), together with a term of \(B(x)\)'' \\
Generalizes & \(A \to B\) (function type) & \(A \times B\) (product type) \\
Constant-family collapse (\(B\) ignores \(x\)) & ordinary function type \(A \to B\) & ordinary product \(A \times B\) \\
Categorical reading & indexed \textbf{product} & indexed \textbf{coproduct} (disjoint union) \\
Lean surface syntax & \texttt{(x : A) → B x} & \texttt{Σ x : A, B x}, built with \texttt{⟨\_\allowbreak{}, \_\allowbreak{}⟩} \\
Logic special case (\(B : A \to \mathrm{Prop}\)) & \texttt{∀ x, P x} & \texttt{∃ x, P x} \\
Witness/value extraction & always. Applying a Π-typed function to an argument is ordinary evaluation & allowed for \texttt{Sigma} (\texttt{Type}-valued, \texttt{.fst}/\texttt{.snd}); \textbf{not} allowed for \texttt{Exists} (\texttt{Prop}-valued, proof irrelevance forbids it) \\
First example this book gave & \texttt{Vec.replicate : (n : Nat) → Vec α n} (Section 3 above) & the fields of \texttt{Fin n} itself, \texttt{val : Nat} paired with \texttt{isLt : val < n} (above) \\
\end{longtable}

The one asymmetry worth remembering from the row above is that the logic special case (\texttt{∀}) of the Π-type and that of the Σ-type (\texttt{∃}) both extract cleanly as \emph{propositions}, but only the \emph{data} special case (\texttt{Sigma}) of the Σ-type supports extracting its witness back out. \texttt{∃} looks like a Σ-type and reads like one, but proof irrelevance makes it a strictly weaker, non-extractable cousin. The ``caveat about \texttt{∃}'' above is the one place the two columns of this table are not simply mirror images of each other.

\subsection{Recursors and eliminators, named}\label{sec:02-terminology-and-coc:02-pi-sigma-and-coc:6}

\hyperref[sec:02-terminology-and-coc:01-terminology:1]{Chapter 2, Section 1} promised a formal treatment of the \textbf{recursor} (also called an \textbf{eliminator}) once Π/Σ-types were available. Here it is. For an inductive type like \texttt{Nat}, with constructors \texttt{zero : Nat} and \texttt{succ : Nat → Nat}, the \textbf{recursor} \texttt{Nat.rec} is the single term that makes ``define a function, or prove a statement, by giving one case per constructor'' precise:

\begin{lstlisting}
#check @Nat.rec
-- {motive : Nat (*$\rightarrow$*) Sort u}
--   (*$\rightarrow$*) motive Nat.zero
--   (*$\rightarrow$*) ((n : Nat) (*$\rightarrow$*) motive n (*$\rightarrow$*) motive n.succ)
--   (*$\rightarrow$*) (t : Nat) (*$\rightarrow$*) motive t
\end{lstlisting}

Read the type of \texttt{Nat.rec} itself as a Π-type over a \textbf{motive} \texttt{motive : Nat → Sort u} (the ``motive'' of Section 4, spelled out in full). Supply a ``zero case'' (a term of \texttt{motive Nat.zero}), a ``succ case'' (a function turning a value/proof for \texttt{n} into one for \texttt{n.succ}), and \texttt{Nat.rec} produces a term of \texttt{motive t} for \emph{any} \texttt{t : Nat}. This is structural induction/recursion stated as one Π-typed term rather than as a separate principle bolted onto the language. Proving something for \texttt{zero} and showing it is preserved by \texttt{succ} is \emph{literally} supplying the two remaining arguments of \texttt{Nat.rec}.

\textbf{A first worked example.} The simplest motive is a \emph{constant} one. The result type does not actually depend on which \texttt{Nat} was supplied, only its value does. Doubling a number, written directly as an application of \texttt{Nat.rec} instead of via \texttt{+}:

\begin{lstlisting}
def double (n : Nat) : Nat :=
  Nat.rec 0 (fun _ ih => ih + 2) n

#eval double 5   -- 10
\end{lstlisting}

Here \texttt{motive := fun \_\allowbreak{} => Nat} (inferred, since the zero case \texttt{0} and the result of the succ case both have type \texttt{Nat}). The zero case supplies \texttt{0} for \texttt{double 0}, and the succ case says ``given the doubled value \texttt{ih} for \texttt{n}, the doubled value for \texttt{n.succ} is \texttt{ih + 2},'' exactly the recursive definition of doubling, with no separate ``recursion'' mechanism beyond \texttt{Nat.rec} itself. A \texttt{dbg\_\allowbreak{}trace} inside the succ case makes each of the five hidden recursive steps visible:

\begin{lstlisting}
def double' (n : Nat) : Nat :=
  Nat.rec 0 (fun _ ih => dbg_trace s!"double: succ case, ih={ih}, adding 2"; ih + 2) n

#eval double' 5
-- double: succ case, ih=0, adding 2
-- double: succ case, ih=2, adding 2
-- double: succ case, ih=4, adding 2
-- double: succ case, ih=6, adding 2
-- double: succ case, ih=8, adding 2
-- 10
\end{lstlisting}

Five lines print, one per \texttt{succ} layer \texttt{Nat.rec} peels off \texttt{5 = succ(succ(succ(succ(succ zero))))}, each showing the running total \emph{before} the final \texttt{+ 2}. That total is \texttt{0}, then \texttt{2}, then \texttt{4}, then \texttt{6}, then \texttt{8}, then the returned \texttt{10}. Nothing about \texttt{Nat.rec} was changed to make this possible. The trace prints from inside the ordinary succ-case function supplied as the second argument of \texttt{Nat.rec}, the same function already described above.

\textbf{A second worked example, over a different inductive type.} The same pattern generalizes to any inductive type, not just \texttt{Nat}. Computing the length of a list via \texttt{List.rec} directly, rather than via the built-in \texttt{List.length}:

\begin{lstlisting}
noncomputable def myLength {(*$\alpha$*) : Type} (l : List (*$\alpha$*)) : Nat :=
  List.rec (motive := fun _ => Nat) 0 (fun _ _ tailLen => tailLen + 1) l

#reduce myLength [1, 2, 3]        -- 3
#reduce myLength ([] : List Nat)  -- 0
\end{lstlisting}

\texttt{noncomputable} and \texttt{\#reduce} in place of \texttt{\#eval} here are a Lean implementation detail. The code generator that backs \texttt{\#eval} does not yet support compiling \texttt{List.rec} directly, so this definition is checked and reduced by the kernel via \texttt{\#reduce} instead. The mathematical content is identical to \texttt{double} above, just for \texttt{List} in place of \texttt{Nat}. The motive is again constant (\texttt{fun \_\allowbreak{} => Nat}), the ``nil case'' is \texttt{0}, and the ``cons case'' says ``given the length of the tail, \texttt{tailLen}, the length with one more element prepended is \texttt{tailLen + 1}.''

Unlike \texttt{double} above, the recursion of \texttt{myLength} cannot be watched with \texttt{dbg\_\allowbreak{}trace}. A \texttt{dbg\_\allowbreak{}trace} call placed inside the cons case prints nothing at all under \texttt{\#reduce}, verified directly. The definitional reduction performed by the kernel (what \texttt{\#reduce} invokes) is a pure, side-effect-free rewriting process, not a compiled/interpreted evaluation the way \texttt{\#eval} is, so the print action of \texttt{dbg\_\allowbreak{}trace} never fires. This is the same \texttt{noncomputable} boundary from one paragraph up, seen from a second angle. Whatever cannot be compiled also cannot be traced by any mechanism that relies on running compiled code.

\textbf{A counterexample, getting the case order wrong.} The two case arguments of \texttt{Nat.rec} are positional, not named. Supplying them in the wrong order is a genuine, easy mistake, and Lean rejects it as a type error rather than silently doing the wrong thing:

\begin{lstlisting}
def doubleBad (n : Nat) : Nat :=
  Nat.rec (fun _ ih => ih + 2) 0 n
\end{lstlisting}

\begin{lstlisting}
error: Type mismatch
  fun x ih => ih + 2
has type
  (x : ?m.4) (*$\rightarrow$*) (ih : ?m.13 x) (*$\rightarrow$*) ?m.15 x ih
but is expected to have type
  Nat
\end{lstlisting}

The succ-case function was placed where the zero case belongs. Since the first explicit argument of \texttt{Nat.rec} must have type \texttt{motive Nat.zero} (here, plainly \texttt{Nat}), a function is rejected on the spot. This is the same class of error underlying ``motive is not type correct'' messages elsewhere in this book. The argument positions of \texttt{Nat.rec} encode real structure (which case is which), and getting that structure wrong is caught before anything runs, not discovered later at \texttt{\#eval} time.

An \textbf{eliminator} is the general name for this same pattern applied to \emph{any} inductive type. It is the single term that ``uses'' a value of that type by case-splitting on which constructor built it, with a motive tracking what is being proved or built in each case. \texttt{Nat.rec} above is the eliminator for \texttt{Nat}. The \texttt{Or.elim \{P Q R : Prop\} (h : P ∨ Q) (hpr : P → R) (hqr : Q → R) : R} shown in \hyperref[sec:04-propositions-and-proofs:05-and-or-not:1]{Chapter 4, Section 5} is the eliminator for \texttt{Or}, a case for \texttt{Or.inl} and a case for \texttt{Or.inr}, with the (non-dependent, here) motive fixed to the constant type \texttt{R}. Every \texttt{match}/\texttt{cases}/\texttt{induction} used from here on compiles down to exactly this, an application of the eliminator for the relevant inductive type, built automatically rather than written out by hand.

\subsection{The calculus of constructions, named}\label{sec:02-terminology-and-coc:02-pi-sigma-and-coc:7}

Putting the pieces together, the \textbf{calculus of constructions} (CoC, Coquand--Huet), the core type theory underlying Lean, Coq, and similar systems, is a small formal system built from a variable, a function abstraction, and function application, plus the following.

\begin{enumerate}
\def\labelenumi{\arabic{enumi}.}
\tightlist
\item
  An infinite hierarchy of type universes \(\mathtt{Type}\,0,
  \mathtt{Type}\,1, \ldots\), formalized in \hyperref[sec:06-rigor-check:03-typing-rules-and-safety:1]{Chapter 6, Section 3}.
\item
  Π-types generalizing \(\to\), allowing types to depend on terms (Section 3 of this chapter and the top of this page).
\item
  \texttt{inductive} type declarations, in the specific extension used by Lean, the \textbf{Calculus of Inductive Constructions} (CIC), such as \texttt{Nat}, \texttt{Bool}, \texttt{Vec}/\texttt{Fin} (Section 3), \texttt{Path} (Chapter 12), and every \texttt{structure} written throughout, giving Σ-types and much more (arbitrary recursive data) beyond what bare CoC provides.
\item
  \texttt{Prop} as a distinguished, proof-irrelevant universe (Chapter 4, and above), making Curry--Howard a precise correspondence rather than an analogy.
\end{enumerate}

Every single Lean construct used in this book, \texttt{def}, \texttt{structure}, \texttt{theorem}, \texttt{∀}, \texttt{∃}, and tactics (which merely \emph{build} CIC terms, one piece at a time, without the terms being written by hand), compiles down to a term in exactly this system, checked by the Lean kernel using nothing more than typing rules like the ones sketched here and in Chapter 6, applied mechanically. \hyperref[sec:06-rigor-check:03-typing-rules-and-safety:1]{Chapter 6, Section 3} completes the picture with the typing judgments of the simply typed λ-calculus and its progress/preservation theorems, the formal reason ``well-typed proofs do not go wrong,'' plus the universe-formation rule only sketched by name above.

\subsection{Sources, quoted}\label{sec:02-terminology-and-coc:02-pi-sigma-and-coc:8}

Formal definitions and citations for this section, gathered here for reference (full entries in the \hyperref[sec:root:bibliography:1]{Bibliography}):

\begin{itemize}
\tightlist
\item
  \textbf{Σ-type.} ``Given a type \(A : U\) and a family \(B : A \to U\), the dependent pair type is written as \(\sum_{(x:A)} B(x) : U\) \ldots{} If \(B\) is constant, then the dependent pair type is the ordinary Cartesian product type'' (\cite{HoTT2013}, §1.6 ``Dependent pair types (Σ-types)''). Picture it like this. A labeled suitcase where the label decides what kind of compartment to expect inside, a ``shoes'' tag pairs with a shoe-shaped compartment, an ``electronics'' tag with a different one. When every tag happens to pair with the \emph{same} compartment shape regardless of the label, that's just an ordinary suitcase-and-item pair, the constant-family special case, which is the ordinary product \(A \times B\).
\item
  \textbf{\texttt{Prop}.} ``The type \texttt{Prop} is syntactic sugar for \texttt{Sort 0}, the very bottom of the type hierarchy'' (\cite{TPIL4}, §3.1 ``Propositions as Types'').
\item
  \textbf{Proof irrelevance.} ``If \texttt{p : Prop} is any proposition, the Lean kernel treats any two elements \texttt{t1 t2 : p} as being definitionally equal \ldots{} This is known as proof irrelevance'' (\cite{TPIL4}, §3.1). Picture it like this. A ``verified'' stamp on a passport. Customs doesn't care which officer stamped it or how the check was carried out, only that a valid stamp exists. Any two proofs of the same \texttt{P : Prop} are just as interchangeable.
\item
  \textbf{Recursor / eliminator.} The working statement used by this book, built on the calculus of constructions (\cite{CoquandHuet1988}), is the single Π-typed term (\texttt{Nat.rec} and its analogues) that makes ``one case per constructor'' precise. Every \texttt{match}/\texttt{cases}/\texttt{induction} compiles down to one.
\item
  \textbf{Calculus of constructions (CoC).} The working statement used by this book, after Coquand and Huet (\cite{CoquandHuet1988}, §1 ``The Abstract Syntax of Terms,'' §2.1 ``The Inference System of Constructions''), is a small formal system built from a variable, a function abstraction, and function application, plus an infinite hierarchy of universes and Π-types. The specific extension of Lean with \texttt{inductive} type declarations (giving Σ-types and general recursive data) is the Calculus of Inductive Constructions (CIC), due to a later paper, Thierry Coquand and Christine \textbf{Paulin-Mohring}, ``Inductively Defined Types,'' in \emph{COLOG-88}, Lecture Notes in Computer Science 417, Springer, 1990, not the 1988 paper itself. Note that this is a \emph{different} paper from the one by Pfenning \& Paulin-Mohring, ``Inductively Defined Types in the Calculus of Constructions'' (1989), with which it is easily conflated. The surname is hyphenated in both.
\item
  Pierce (\cite{Pierce2002}) is cited here only by analogy. \emph{Types and Programming Languages} is explicitly \emph{not} a dependently-typed textbook (the preface by Pierce himself says dependent types are ``mentioned only in passing,'' developed no further than the one-paragraph sketch in §30.5 of ``families of types indexed by terms''), so it does not cover eliminators/recursors for inductive types in a dependent setting. What it does cover, relevant here only by analogy, is non-dependent \texttt{case} analysis on sum/variant types (§11.9--11.10) and on \texttt{µ}-recursive types like \texttt{NatList = μX.⟨nil:Unit, cons:\{Nat,X\}⟩} (Ch. 21), plus Church encodings of algebraic datatypes (Ch. 6 untyped, §23.4 typed/System F).
\item
  Martin-Löf (\cite{MartinLof1984}) is the foundational source for dependent Π/Σ types and universes predating CoC, for readers wanting the idea in its original, non-CoC-specific form.
\item
  All Lean code in this section was checked directly against the toolchain pinned in \texttt{lean\_\allowbreak{}project/\allowbreak{}lean-\allowbreak{}toolchain} in this repository rather than only described; the \texttt{\#print Fin} output and both error messages shown are copied from real \texttt{lake env lean} runs.
\end{itemize}
